\chapter[Resultados]{Resultados}
\label{cap:resultados}

Este capítulo reporta os resultados dos cinco modelos descritos no Capítulo~\ref{cap:Metodologia} sobre o benchmark Families in the Wild (FIW) com o split Track-I oficial do RFIW. O Modelo 02 também aparece em KinFaceW-I com validação cruzada de cinco dobras, para situar o trabalho contra a literatura clássica. A unidade experimental reportada é a melhor checkpoint salva em validação por run, e a métrica principal é o ROC AUC no conjunto de teste, complementada por Average Precision, TAR@FAR=0,001, TAR@FAR=0,01, TAR@FAR=0,1 e accuracy por tipo de relação. Os baselines com modelos multimodais entram em seção própria com tarefa diferente (classificação em onze classes), e os números não são comparados linha a linha com os supervisionados.

O conjunto inclui um baseline mínimo sem treinamento (B0, AdaFace IR-101 congelado com cosseno), o Modelo 02 (ViT com cross-attention bidirecional FaCoR), o Modelo 05 (DINOv2 com LoRA e cross-attention diferencial), o Modelo 06 (retrieval-augmented sobre galeria de pares positivos) e o Modelo 12 (RGCK-Net, com AdaFace IR-101 parcialmente descongelado e tokens regionais anatômicos). A escolha desses cinco cobre o piso sem treino, dois caminhos paramétricos plenos (M02 full FT, M12 partial FT), um caminho parameter-efficient (M05) e um caminho não-paramétrico (M06).

\section{Visão geral}

A Tabela~\ref{tab:overview_runs} resume o melhor run de cada modelo em FIW Track-I, ordenado por Test ROC AUC. A linha do Modelo 12 aparece em destaque porque corresponde ao maior AUC observado neste conjunto experimental, mas com a ressalva de que a margem sobre o Modelo 02 é incremental (+0,006) e não foi validada por múltiplas seeds ou k-fold em FIW.

\begin{table}[h]
\centering
\caption{Visão geral dos cinco modelos avaliados em FIW Track-I, ordenados por Test ROC AUC.}
\label{tab:overview_runs}
\begin{tabular}{lllrrr}
\hline
\textbf{Modelo} & \textbf{Família} & \textbf{Melhor run} & \textbf{Test AUC} & \textbf{TAR@FAR=0,01} & \textbf{Params treináveis} \\
\hline
M12 & RGCK-Net AdaFace + tokens regionais & R002 & \textbf{0{,}856} & \textbf{0{,}176} & 31{,}6M \\
M02 & ViT + FaCoR cross-attention         & R031 & 0{,}850          & 0{,}130          & {\raise.17ex\hbox{$\scriptstyle\sim$}}86M \\
M05 & DINOv2 + LoRA + DiffAttn            & R005 & 0{,}822          & 0{,}100          & 93{,}5M \\
B0  & AdaFace Cosine Frozen               & ---  & 0{,}799          & 0{,}218          & 0 \\
M06 & Retrieval-augmented (ViT frozen)    & R001 & 0{,}776          & 0{,}062          & 8{,}16M \\
\hline
\end{tabular}
\end{table}

O baseline sem treinamento B0 atinge 0,799 de AUC apenas com cosseno sobre embeddings AdaFace pré-treinados, o que situa o piso útil do problema acima de boa parte das tentativas treinadas registradas no projeto. O ganho do Modelo 12 sobre B0 é de +0,057 AUC (cerca de 7\% relativo), enquanto o ganho sobre o Modelo 02 é de +0,006. Ambas as diferenças são reais nas métricas independentes de threshold, mas a segunda é pequena e deve ser lida como melhor resultado observado, não como superioridade definitiva.

\section{Resultados em KinFaceW-I}

KinFaceW-I serve neste trabalho como banco de prototipagem e como referência à literatura clássica de verificação de parentesco. O Modelo 02 foi avaliado no split padrão de cinco dobras e em uma validação cruzada de cinco dobras independente para medir variabilidade entre folds. A Tabela~\ref{tab:kinface_resultados} reporta as duas leituras.

\begin{table}[h]
\centering
\caption{Resultados do Modelo 02 em KinFaceW-I. A primeira coluna usa o checkpoint do Run 023 sobre o split padrão; a segunda reporta média mais ou menos desvio padrão sobre cinco dobras independentes do Run 026.}
\label{tab:kinface_resultados}
\begin{tabular}{lrr}
\hline
\textbf{Métrica} & \textbf{Modelo 02 R023 (split padrão)} & \textbf{Modelo 02 R026 (5-fold CV)} \\
\hline
Accuracy   & 74{,}1\% & 73{,}6\% $\pm$ 5{,}8\% \\
F1         & 0{,}753  & 0{,}771 $\pm$ 0{,}028 \\
ROC AUC    & 0{,}861  & 0{,}842 $\pm$ 0{,}012 \\
Precision  & 71{,}9\% & 69{,}4\% $\pm$ 6{,}3\% \\
Recall     & 79{,}0\% & 87{,}8\% $\pm$ 6{,}2\% \\
Threshold  & 0{,}900  & 0{,}890 $\pm$ 0{,}020 \\
\hline
\end{tabular}
\end{table}

O desvio padrão do AUC sobre as cinco dobras fica em 0,012, e a média entre folds (0,842) é compatível com o valor do split padrão (0,861) dentro dessa variabilidade. Esse resultado dá um piso de confiança na abordagem ViT mais cross-attention para o domínio, e serve como ponto de comparação com a literatura, que historicamente reporta números próximos de 73\% a 80\% de accuracy em KinFaceW-I.

\section{Resultados em FIW Track-I}

FIW Track-I é o benchmark principal deste trabalho. A Tabela~\ref{tab:fiw_main} consolida métricas independentes de threshold (AUC, Average Precision, três níveis de TAR@FAR) para os cinco modelos. As métricas dependentes de threshold ficam na seção sobre operating points, porque comparações diretas entre M02 com $\tau = 0{,}90$ e M12 com $\tau = 0{,}50$ misturam efeito de método e efeito de operating point.

\begin{table}[h]
\centering
\caption{Resultados em FIW Track-I. Métricas independentes de threshold permitem comparação direta entre métodos.}
\label{tab:fiw_main}
\begin{tabular}{lrrrrr}
\hline
\textbf{Modelo} & \textbf{Test AUC} & \textbf{Avg Prec.} & \textbf{TAR@FAR=0,001} & \textbf{TAR@FAR=0,01} & \textbf{TAR@FAR=0,1} \\
\hline
M12 R002 & \textbf{0{,}856} & \textbf{0{,}839} & 0{,}042          & \textbf{0{,}176} & \textbf{0{,}571} \\
M02 R031 & 0{,}850          & 0{,}817          & 0{,}025          & 0{,}140          & 0{,}499 \\
M05 R001 & 0{,}806          & ---              & ---              & \textbf{0{,}152} & --- \\
B0       & 0{,}799          & 0{,}809          & \textbf{0{,}071} & \textbf{0{,}218} & 0{,}524 \\
M06 R001 & 0{,}776          & ---              & ---              & 0{,}062          & --- \\
\hline
\end{tabular}
\end{table}

Três leituras saem da tabela. O Modelo 12 lidera em AUC, Average Precision, TAR@FAR=0,01 (entre os treinados) e TAR@FAR=0,1, e essa é a base para tratá-lo como melhor resultado observado no conjunto experimental. O Modelo 02 segue muito próximo em AUC, com diferença de 0,006, ainda dentro do que se esperaria de variabilidade de single-run em FIW Track-I. O baseline sem treino B0 tem o maior TAR@FAR=0,001 de todos os modelos (0,071), o que reflete o que AdaFace foi treinado para fazer originalmente, separação de identidades em regimes de baixa taxa de falsos positivos, e expõe que parte expressiva do sinal de parentesco já está codificada no embedding de reconhecimento facial sem treino específico.

Comparações diretas entre M12 e M02 em métricas dependentes de threshold são delicadas. O Modelo 12 R002 foi testado com threshold 0,500, padrão do script, porque o metadado de validação não foi gravado no checkpoint antes do treino ser interrompido manualmente. O Modelo 02 R031 opera em threshold 0,900, escolhido em validação para maximizar F1. Isso significa que precision e recall em pontos distintos da curva. No threshold 0,500 do M12 a precision é 79,8\% e o recall é 69,0\%; no threshold 0,900 do M02 a precision é 66,5\% e o recall é 94,1\%. As métricas dependentes de threshold deste trabalho são reportadas com essa ressalva.

\section{Estudo de regime no Modelo 05}

O Modelo 05 passou por sete runs cobrindo quatro regimes de plasticidade do backbone DINOv2 e quatro variações de loss e amostragem. A Tabela~\ref{tab:m05_runs} resume cada run, isolando uma hipótese por vez.

\begin{table}[h]
\centering
\caption{Sete runs do Modelo 05 em FIW Track-I, cada um isolando uma hipótese sobre o backbone, a loss ou o sampler.}
\label{tab:m05_runs}
\begin{tabular}{lp{6.0cm}rrr}
\hline
\textbf{Run} & \textbf{Hipótese isolada} & \textbf{Test AUC} & \textbf{TAR@FAR=0,01} & \textbf{Params} \\
\hline
R001 & LoRA rank 8, backbone congelado, dropout 0     & 0{,}806 & \textbf{0{,}152} & 8{,}47M \\
R002 & LoRA dropout 0,1 + relation weight 0,4         & 0{,}799 & 0{,}095 & 8{,}47M \\
R003 & Class-balanced sampler                          & 0{,}809 & 0{,}098 & 8{,}47M \\
R004 & Partial unfreeze dos 4 blocos finais           & 0{,}812 & 0{,}119 & 36{,}8M \\
R005 & Full unfreeze com LR 5e-6 estilo M02           & \textbf{0{,}822} & 0{,}100 & 93{,}5M \\
R006 & Heavy contrastive loss (peso 0,9)              & 0{,}814 & 0{,}094 & 8{,}47M \\
R007 & Backbone híbrido DINOv2 + ViT do M02           & 0{,}810 & 0{,}136 & 11{,}58M \\
\hline
\end{tabular}
\end{table}

Há um teto arquitetônico em torno de 0,81-0,82 AUC que sete intervenções diferentes não conseguem mover de forma consistente. A R005, com full unfreeze do DINOv2 e LR 5e-6, atinge o maior AUC do M05 mas perde o melhor TAR@FAR=0,01 para a R001, que mantém o regime original LoRA-only. O checkpoint mais útil do Modelo 05 não é o de maior AUC, é a R001, porque o ganho de R005 em AUC vem acompanhado de perda em precisão de alta exigência (0,152 cai para 0,100). A R007 testa a hipótese mais ambiciosa, fundir o DINOv2 com o ViT do Modelo 02 R031 em arquitetura híbrida com cross-attention diferencial intra-face, e produz AUC idêntico ao R001 sem o backbone adicional. Esse achado negativo, descrito como hipótese H0 confirmada no log de execução, indica que as features kinship-fortes do M02 R031 não estão isoladas no ViT mas vêm do sistema completo treinado em conjunto.

\section{Retrieval-augmented no Modelo 06}

O Modelo 06 testa se parte do sinal necessário para verificação de parentesco pode vir de memória não-paramétrica em vez de pesos. A Run 001 mantém o encoder ViT-B/16 congelado e usa uma galeria de 33.207 pares positivos do split de treino do FIW, recuperando 32 pares mais similares por consulta de inferência. A Tabela~\ref{tab:m06_per_rel} mostra a accuracy por relação, que é o ponto mais distintivo do modelo.

\begin{table}[h]
\centering
\caption{Accuracy do Modelo 06 R001 por tipo de relação em FIW Track-I.}
\label{tab:m06_per_rel}
\begin{tabular}{lrr}
\hline
\textbf{Relação} & \textbf{Accuracy} & \textbf{N pares} \\
\hline
sibs (irmãos misto)  & 87{,}2\% & 234 \\
bb (irmão-irmão)     & 86{,}5\% & 860 \\
ss (irmã-irmã)       & 86{,}3\% & 731 \\
md (mãe-filha)       & 85{,}9\% & 1.038 \\
ms (mãe-filho)       & 83{,}9\% & 1.036 \\
gfgs (avô-neto)      & 82{,}7\% & 98 \\
fs (pai-filho)       & 78{,}8\% & 1.135 \\
fd (pai-filha)       & 76{,}9\% & 918 \\
gfgd (avô-neta)      & 75{,}4\% & 138 \\
gmgd (avó-neta)      & 63{,}4\% & 123 \\
gmgs (avó-neto)      & 61{,}2\% & 121 \\
\hline
\end{tabular}
\end{table}

A faixa entre 61,2\% e 87,2\% é a mais estreita entre os modelos com per-relação registrada no mesmo threshold. O Modelo 05 R001, no mesmo threshold, cai a 39,8\% em gfgs como pior classe, deixando um piso muito mais baixo para relações raras. A interpretação é que retrieval-augmentation atua como regularizador estrutural para classes com poucos exemplos. Expor explicitamente pares positivos similares no momento da inferência reduz a pressão sobre os pesos para codificar todas as variações de cada relação. A Run 002 do Modelo 06 testou substituir o backbone por DINOv2 e dobrar K para 64, mas o AUC caiu de 0,776 para 0,731 (-0,045) e o gap val-teste cresceu de 0,060 para 0,090. Backbone mais rico e maior contexto não compensaram a fragilidade da cross-attention contra a galeria fixa.

\section{RGCK-Net no Modelo 12}

O Modelo 12 implementa a proposta RGCK-Net (Region-Guided Cross Kinship Network) descrita no Capítulo~\ref{cap:Metodologia}. A avaliação ocorreu em duas fases consecutivas, uma com o backbone AdaFace IR-101 totalmente congelado (R001) e outra com partial unfreeze do último estágio e da camada de saída (R002). A Tabela~\ref{tab:m12_runs} compara as duas fases.

\begin{table}[h]
\centering
\caption{Duas runs do Modelo 12 em FIW Track-I, isolando o efeito do partial unfreeze do último estágio do AdaFace.}
\label{tab:m12_runs}
\begin{tabular}{lp{5.5cm}rrr}
\hline
\textbf{Run} & \textbf{Regime} & \textbf{Val AUC} & \textbf{Test AUC} & \textbf{Params treináveis} \\
\hline
R001 & Backbone totalmente congelado, treino apenas em tokens, cross-attention, gate e cabeça & 0{,}835 & 0{,}746 & 5{,}59M (7,9\%) \\
R002 & Partial unfreeze: stages 1-3 congelados, stage 4 e output\_layer treináveis, LR 1e-5 & \textbf{0{,}932} & \textbf{0{,}856} & 31{,}55M (44,6\%) \\
\hline
\end{tabular}
\end{table}

A R001 produz val→teste gap de -0,089, o menor entre os modelos AdaFace-based do projeto, mas o teto de validação ficou em 0,835 e o Test AUC em 0,746, abaixo dos modelos paramétricos plenos. A R002 destrava os 26 milhões de parâmetros adicionais correspondentes ao stage 4 mais o output\_layer e reduz o LR de 1e-4 para 1e-5. O Val AUC sobe para 0,932, o maior do projeto, e o Test AUC sobe para 0,856.

A hipótese registrada antes do experimento era de que partial unfreeze elevaria o teto de validação preservando boa parte do gap pequeno do R001. O resultado observado confirmou a direção, embora a magnitude tenha sido maior do que o esperado. O LR conservador de 1e-5 protege as features identitárias profundas do AdaFace de serem reescritas drasticamente, enquanto o sinal de gradiente é forte o suficiente para que o stage 4 e o output\_layer se ajustem ao domínio de parentesco. A redação não atribui o ganho exclusivamente ao LR porque não houve ablação isolada de LR sobre o mesmo conjunto de parâmetros descongelados; o registro experimental sustenta a interação entre partial unfreeze e LR conservador, não a contribuição individual de cada um.

A estrutura regional, com 5 tokens anatômicos (global, olhos, nariz, boca, mandíbula) e cross-region attention 5x5, fornece um prior estrutural que, ao contrário das 196 patch tokens do M02, restringe explicitamente onde o modelo deve comparar as duas faces. Em validação, todas as 11 classes ficam entre 86\% e 96\% de accuracy na época do checkpoint, faixa mais estreita do que qualquer outro modelo do projeto. Esse comportamento sustenta a hipótese da proposta de que regiões anatômicas fornecem indutivo útil para parentesco, embora os números de teste no threshold 0,500 não permitam afirmar isso categoricamente porque o threshold não foi otimizado em validação.

\section{Accuracy por relação}

A Tabela~\ref{tab:per_rel_compare} compara accuracy por tipo de relação entre modelos no threshold 0,500. M02 R031 não aparece nesta tabela porque opera em threshold 0,900 e os números por classe não são comparáveis diretamente; aparece em texto como referência ao final.

\begin{table}[h]
\centering
\caption{Accuracy por tipo de relação em FIW Track-I no threshold 0,500, em pontos percentuais. As quatro classes inferiores são relações de segundo grau (avós e netos).}
\label{tab:per_rel_compare}
\begin{tabular}{lrrrr}
\hline
\textbf{Relação} & \textbf{N pares} & \textbf{B0} & \textbf{M05 R001} & \textbf{M12 R002} \\
\hline
fs (pai-filho)      & 1.135 & 89{,}9\% & 71{,}6\% & 68{,}6\% \\
fd (pai-filha)      & 918   & 80{,}1\% & 71{,}5\% & 68{,}4\% \\
ms (mãe-filho)      & 1.036 & 86{,}6\% & 69{,}4\% & 69{,}4\% \\
md (mãe-filha)      & 1.038 & 89{,}4\% & 67{,}3\% & 68{,}7\% \\
bb (irmão-irmão)    & 860   & 89{,}9\% & 79{,}8\% & 75{,}4\% \\
ss (irmã-irmã)      & 731   & 89{,}9\% & 77{,}2\% & 75{,}9\% \\
sibs (irmãos misto) & 234   & 90{,}6\% & 83{,}3\% & 79{,}5\% \\
\hline
gfgs (avô-neto)     & 98  & 68{,}4\% & 39{,}8\% & 39{,}8\% \\
gfgd (avô-neta)     & 138 & 72{,}5\% & 50{,}7\% & \textbf{52{,}2\%} \\
gmgs (avó-neto)     & 121 & 52{,}9\% & 52{,}1\% & 44{,}6\% \\
gmgd (avó-neta)     & 123 & 75{,}6\% & 40{,}7\% & 36{,}6\% \\
\hline
non-kin             & 6.993 & 48{,}6\% & --- & 84{,}0\% \\
\hline
\end{tabular}
\end{table}

A leitura cruzada precisa de cuidado. No threshold 0,500, B0 opera em regime de alta recall e baixa precision; o modelo classifica como kin a maioria dos pares positivos (53\% a 91\% de accuracy em classes kin) mas erra em mais da metade dos pares non-kin (accuracy non-kin de 48,6\%). Esse comportamento é típico de cosseno sobre embeddings AdaFace, em que pares de irmãos e de pais e filhos têm similaridade alta independente do parentesco real. M12 R002 troca recall por precisão: a accuracy em classes kin cai (68\% a 80\%) mas a accuracy em non-kin sobe para 84\%, o que sustenta o ganho em AUC observado anteriormente. A relação gfgd, historicamente uma das mais difíceis do projeto, vai de 31\% no M09 R001 (modelo AdaFace anterior) a 52\% no M12 R002. As demais relações de avós e netos seguem como o principal gargalo, com gmgd em 36,6\% e gmgs em 44,6\% no M12, ainda longe de qualquer paridade entre classes.

\section{Baselines com modelos multimodais}

Dois VLMs avaliam o que um modelo multimodal generalista consegue inferir sobre parentesco facial sem treinamento específico. A tarefa é classificação em onze classes do FIW, não verificação binária, e por isso os números não são comparados diretamente com os modelos supervisionados. A Tabela~\ref{tab:vlm_resultados} resume.

\begin{table}[h]
\centering
\caption{Resultados dos baselines com VLMs em classificação multiclasse de onze relações.}
\label{tab:vlm_resultados}
\begin{tabular}{lllr}
\hline
\textbf{Modelo} & \textbf{Regime} & \textbf{Amostra} & \textbf{Accuracy} \\
\hline
Codex gpt-5.4-mini & Zero-shot                          & 750 pares balanceados & 33{,}1\% \\
Codex gpt-5.4-mini & Adaptação por prompt (7-shot)      & 750 pares (mesmos)    & 26{,}7\% \\
Codex gpt-5.4-mini & Prompts direcionados por treino    & 750 pares (mesmos)    & 28{,}4\% \\
Claude Sonnet 4.6  & Zero-shot direto                   & 75 pares (piloto)     & 37{,}3\% \\
\hline
\end{tabular}
\end{table}

A accuracy do Codex em zero-shot é 33,1\%, e o experimento de adaptação ao domínio via few-shot mais prompt estruturado mais calibração leve piorou o desempenho para 26,7\% nas mesmas amostras. As três configurações de prompt testadas mantiveram a queda, e o relatório técnico associado interpreta o efeito como amplificação de heurísticas de idade e gênero quando o modelo é forçado a produzir saídas auxiliares antes da resposta final. As quatro classes gfgs, gfgd, gmgs e gmgd ficam em 0\% nos dois modelos, em todas as configurações testadas, indicando que a aparência facial isolada não basta para distinções genealógicas de segundo grau em zero-shot. O Claude Sonnet só foi avaliado em piloto de 75 pares, então o número de 37,3\% deve ser lido como indicação, não como comparação robusta.

\section{Síntese dos resultados}

O conjunto experimental sustenta cinco leituras principais. O Modelo 12 obteve o maior Test AUC observado neste estudo (0,856) e lidera Average Precision e TAR@FAR=0,01 entre os modelos treinados, mas com margem incremental sobre o Modelo 02 (+0,006 AUC) que requer confirmação por múltiplas seeds ou validação cruzada em FIW antes de qualquer afirmação categórica de superioridade. O Modelo 02 ViT full FT segue competitivo (0,850) e mantém o operating point que historicamente maximiza F1 em verificação binária. O baseline sem treinamento B0 atinge AUC 0,799 e o maior TAR@FAR=0,001 do projeto, o que mostra que o AdaFace pré-treinado já codifica parte expressiva do sinal de parentesco e que o ganho arquitetônico do M12 sobre esse piso é de +0,057, leitura mais informativa do que o ganho sobre o M02. O Modelo 05 com LoRA é Pareto-relevante em eficiência paramétrica, atingindo TAR@FAR=0,01 = 0,152 com 8,47 milhões de parâmetros treináveis. O Modelo 06 retrieval-augmented entrega o perfil per-relação mais uniforme entre os métodos avaliados, e os VLMs zero-shot ficam claramente abaixo dos modelos treinados e falham por completo em relações de avós e netos.

A literatura registrada sobre o FIW reporta resultados próximos de 0,85 a 0,92 em accuracy ou AUC dependendo do protocolo. Os resultados deste trabalho situam-se nessa faixa em AUC com M12 e M02, e abaixo dela quando se considera accuracy no operating point 0,500. As métricas independentes de threshold são o referencial mais defensável para a comparação entre métodos deste estudo, e accuracy ou F1 só devem ser interpretadas dentro do mesmo operating point.
